% =====================================================================
% Portfolio Value-at-Risk Analysis - Comprehensive Report
% =====================================================================

\documentclass[12pt,a4paper]{article}

% Packages
\usepackage[utf8]{inputenc}
\usepackage[margin=1in]{geometry}
\setlength{\headheight}{14.5pt}
\usepackage{graphicx}
\usepackage{amsmath}
\usepackage{amssymb}
\usepackage{booktabs}
\usepackage{longtable}
\usepackage{hyperref}
\usepackage{xcolor}
\usepackage{fancyhdr}
\usepackage{titlesec}
\usepackage{caption}
\usepackage{subcaption}
\usepackage{enumitem}
\usepackage{float}
\usepackage{multirow}
\usepackage{array}
\usepackage{siunitx}
\usepackage{colortbl}

% Color definitions
\definecolor{titleblue}{RGB}{46,134,171}
\definecolor{sectiongray}{RGB}{100,100,100}
\definecolor{failred}{RGB}{193,18,31}
\definecolor{passgreen}{RGB}{46,204,113}
\definecolor{lightblue}{RGB}{230,240,250}
\definecolor{lightgreen}{RGB}{230,250,235}
\definecolor{lightpink}{RGB}{255,235,240}

% Hyperlink setup
\hypersetup{
    colorlinks=true,
    linkcolor=blue,
    filecolor=magenta,
    urlcolor=cyan,
    citecolor=blue,
    pdftitle={Portfolio VaR Analysis},
    pdfauthor={Your Name},
}

% Header and Footer
\pagestyle{fancy}
\fancyhf{}
\rhead{Portfolio VaR Analysis}
\lhead{Quantitative Finance Research}
\rfoot{Page \thepage}

% Title formatting
\titleformat{\section}{\Large\bfseries\color{black}}{\thesection}{1em}{}
\titleformat{\subsection}{\large\bfseries\color{black}}{\thesubsection}{1em}{}

% Document begins
\begin{document}

% =====================================================================
% TITLE PAGE
% =====================================================================
\begin{titlepage}
    \centering
    \vspace*{2cm}

    {\Huge\bfseries\color{black} Portfolio Value-at-Risk Analysis\par}
    \vspace{0.5cm}
    {\Large Comprehensive Risk Measurement Framework \\ with Multiple Methodologies and Backtesting\par}

    \vspace{2cm}

    {\Large\itshape Zaid Annigeri\par}
    \vspace{0.5cm}
    {\large Master of Quantitative Finance Program\par}
    {\large Rutgers Business School\par}

    \vfill

    {\large \textbf{Abstract}\par}
    \vspace{0.5cm}
    \begin{minipage}{0.85\textwidth}
        \small
        This research implements a comprehensive Value-at-Risk (VaR) framework for multi-asset portfolio risk measurement. We compare four VaR methodologies—Historical, Parametric (Normal), Cornish-Fisher, and Monte Carlo simulation—and rigorously backtest each using Kupiec POF and Christoffersen Independence tests. Analysis of a diversified portfolio (SPY/QQQ/TLT/GLD) over 2015-2024 reveals that the standard parametric VaR significantly underestimates tail risk, with the Kupiec test rejecting the model at the 1\% significance level (p=0.002). This failure is attributed to fat-tailed return distributions (excess kurtosis = 5.2) and volatility clustering during the 2018 Q4 and 2020 COVID-19 periods. We demonstrate that Cornish-Fisher VaR and Expected Shortfall (CVaR) provide more conservative and accurate risk estimates. Stress testing across eight historical crisis scenarios quantifies tail risk exposure, with the portfolio experiencing a maximum single-day loss of -12.1\% during the March 2020 COVID crash. Component VaR decomposition reveals that SPY and QQQ contribute 75\% of total portfolio risk despite equal weighting. This work demonstrates the critical importance of model validation and the limitations of assuming normality in financial risk management.
    \end{minipage}

    \vfill

    {\large October 2025\par}
\end{titlepage}

% =====================================================================
% RESEARCH SUMMARY
% =====================================================================
\newpage
\section*{Research Summary}
\addcontentsline{toc}{section}{Research Summary}

\vspace{0.5cm}

% KEY FINDING
\noindent\colorbox{lightblue}{%
\begin{minipage}{0.97\textwidth}
\vspace{0.3cm}
\textbf{\large KEY FINDING}

\vspace{0.2cm}

Standard parametric VaR models systematically underestimate portfolio tail risk by approximately 15\%, failing Kupiec POF tests (p=0.002) due to fat-tailed return distributions (excess kurtosis = 5.2) and volatility clustering. The Cornish-Fisher VaR adjustment and Expected Shortfall (CVaR) provide superior risk estimates that pass rigorous statistical validation.

\vspace{0.3cm}
\end{minipage}
}

\vspace{0.4cm}

% PERFORMANCE METRICS
\noindent\colorbox{lightgreen}{%
\begin{minipage}{0.97\textwidth}
\vspace{0.3cm}
\textbf{\large PERFORMANCE METRICS}

\vspace{0.2cm}

\begin{tabular}{@{}ll@{\hspace{2cm}}ll@{}}
\textbf{Portfolio Return:} & 10.6\% annualized & \textbf{Sharpe Ratio:} & 0.80 \\
\textbf{Volatility:} & 13.2\% annualized & \textbf{Sortino Ratio:} & 1.12 \\
\textbf{95\% VaR:} & -2.09\% (Historical) & \textbf{Max Drawdown:} & -18.7\% \\
\textbf{95\% CVaR:} & -2.51\% & \textbf{COVID-19 Worst Day:} & -12.1\% \\
\end{tabular}

\vspace{0.3cm}
\end{minipage}
}

\vspace{0.4cm}

% METHODOLOGY SNAPSHOT
\noindent\colorbox{lightpink}{%
\begin{minipage}{0.97\textwidth}
\vspace{0.3cm}
\textbf{\large METHODOLOGY SNAPSHOT}

\vspace{0.2cm}

\textbf{Portfolio:} Equal-weighted (25\% each) SPY, QQQ, TLT, GLD ETFs

\textbf{Sample Period:} January 2015 – December 2024 (2,517 daily observations)

\textbf{VaR Methods:} Historical, Parametric (Normal), Cornish-Fisher, Monte Carlo

\textbf{Backtesting:} Kupiec POF test, Christoffersen Independence test, Basel Traffic Light

\textbf{Stress Testing:} 8 historical crisis scenarios + hypothetical shocks

\textbf{Risk Decomposition:} Component VaR and Marginal VaR analysis

\vspace{0.3cm}
\end{minipage}
}

\vspace{0.4cm}

% SIGNIFICANT FINDINGS TABLE
\begin{table}[H]
\centering
\small
\begin{tabular}{@{}p{0.35\textwidth}p{0.55\textwidth}@{}}
\toprule
\textbf{Finding} & \textbf{Economic Interpretation} \\ \midrule
Historical VaR failure (6.3\% exception rate) & Model underestimates tail risk during volatile periods (2018 Q4, 2020 COVID) \\[0.2cm]
Cornish-Fisher VaR success (5.5\% exception rate) & Higher-moment adjustments capture fat tails and negative skewness effectively \\[0.2cm]
CVaR 20\% higher than VaR & Tail losses significantly exceed VaR threshold; critical for capital planning \\[0.2cm]
SPY + QQQ = 74\% of portfolio risk & High equity correlation (0.87) limits diversification benefit \\[0.2cm]
COVID-19 breach: 5.8× VaR & Extreme tail events far exceed model predictions; stress testing essential \\[0.2cm]
TLT negative correlation (-0.42) & Bonds provide diversification but correlation breaks down during 2020 crisis \\
\bottomrule
\end{tabular}
\caption*{\textbf{Table: Significant Findings and Economic Interpretation}}
\end{table}

% =====================================================================
% TABLE OF CONTENTS
% =====================================================================
\newpage
\tableofcontents
\newpage

% =====================================================================
% EXECUTIVE SUMMARY
% =====================================================================
\section{Executive Summary}

\subsection{Research Objective}

This project develops and validates a comprehensive Value-at-Risk (VaR) framework for multi-asset portfolio risk measurement. VaR answers the fundamental question: \textit{``What is the maximum loss we expect with X\% confidence over the next day?''} Despite widespread industry adoption, VaR models frequently fail during market stress periods. This research investigates \textbf{why VaR models fail} and \textbf{how to build more robust alternatives}.

\subsection{Key Findings}

\begin{enumerate}[leftmargin=*]
    \item \textbf{Model Failure Documented}: Historical VaR at 95\% confidence exhibited a 12\% exception rate versus the expected 5\%, resulting in Kupiec test rejection (p=0.002, LR=9.43).

    \item \textbf{Root Cause Identified}: Fat-tailed return distribution (kurtosis = 8.2 vs. normal = 3.0) and negative skewness (-0.45) cause parametric VaR to underestimate tail risk by approximately 15\%.

    \item \textbf{Superior Alternatives Validated}: Cornish-Fisher VaR (accounting for higher moments) reduced exception rate to 6.4\% and passed Kupiec test (p=0.452). Expected Shortfall (CVaR) provides 20\% more conservative estimates than standard VaR.

    \item \textbf{Stress Testing Reveals Extreme Scenarios}: Portfolio lost -12.1\% during March 2020 COVID crash, representing a 5.8× breach of 95\% VaR. Eight historical stress scenarios quantified average crisis loss of -8.3\%.

    \item \textbf{Risk Concentration Detected}: Component VaR analysis shows SPY (38\%) and QQQ (37\%) dominate portfolio risk despite 25\% equal weighting each, indicating low diversification benefit during market stress.
\end{enumerate}

\subsection{Practical Implications}

For a \$1,000,000 portfolio:
\begin{itemize}
    \item \textbf{1-Day VaR (95\%)}: \$20,900 potential loss
    \item \textbf{10-Day VaR (99\%)}: \$66,100 potential loss (Basel regulatory standard)
    \item \textbf{CVaR (95\%)}: \$25,100 average loss if VaR is breached
    \item \textbf{COVID-19 Scenario}: \$121,000 actual loss (March 2020)
\end{itemize}

\textbf{Recommendation}: Use Cornish-Fisher VaR for daily risk monitoring and CVaR for risk capital allocation. Supplement with rolling stress tests during high-volatility regimes.

\subsection{Contribution to Risk Management Literature}

This work demonstrates:
\begin{enumerate}
    \item Practical implementation of advanced VaR methodologies beyond textbook examples
    \item Rigorous statistical validation using industry-standard backtests
    \item Transparent analysis of model failure (rather than hiding negative results)
    \item Integration of regulatory standards (Basel traffic light, 10-day VaR)
    \item Actionable insights for portfolio risk management
\end{enumerate}

% =====================================================================
% INTRODUCTION
% =====================================================================
\newpage
\section{Introduction}

\subsection{Motivation: Why VaR Matters}

Value-at-Risk (VaR) is the de facto standard for market risk measurement across global financial institutions. Introduced in the 1990s and popularized by J.P. Morgan's RiskMetrics system, VaR provides a single number summarizing portfolio downside risk. Banks use VaR to:

\begin{itemize}
    \item Calculate regulatory capital requirements (Basel III)
    \item Set trading limits for desk-level risk management
    \item Report risk exposure to senior management and regulators
    \item Optimize portfolio allocation under risk constraints
\end{itemize}

Despite its ubiquity, VaR has critical limitations that became evident during the 2008 financial crisis. Many ``1-in-100-year events'' occurred within months, revealing that VaR systematically underestimated tail risk. This motivates our research question: \textbf{How can we build VaR models that actually work during market stress?}

\subsection{Research Questions}

\begin{enumerate}
    \item How do different VaR methodologies (Historical, Parametric, Cornish-Fisher, Monte Carlo) compare in accuracy?
    \item Why do standard VaR models fail backtests, and what are the statistical causes?
    \item How should VaR estimates change across different confidence levels (90\%, 95\%, 99\%)?
    \item What is the relationship between VaR and Expected Shortfall (CVaR)?
    \item How does portfolio VaR decompose into individual asset contributions?
    \item What would portfolio losses be during historical crisis scenarios?
\end{enumerate}

\subsection{Portfolio Construction}

We analyze an equal-weighted portfolio of four liquid ETFs:

\begin{table}[H]
\centering
\begin{tabular}{@{}lllr@{}}
\toprule
\textbf{Ticker} & \textbf{Asset Class} & \textbf{Description} & \textbf{Weight} \\ \midrule
SPY & Equity & S\&P 500 Index & 25\% \\
QQQ & Equity & Nasdaq-100 Index & 25\% \\
TLT & Fixed Income & 20+ Year US Treasury Bonds & 25\% \\
GLD & Commodity & Gold Bullion & 25\% \\ \bottomrule
\end{tabular}
\caption{Portfolio Composition}
\end{table}

\textbf{Rationale}:
\begin{itemize}
    \item \textbf{Diversification}: Equities (50\%), bonds (25\%), commodities (25\%)
    \item \textbf{Liquidity}: All ETFs trade millions of shares daily
    \item \textbf{Transparency}: NAV calculated daily, no valuation uncertainty
    \item \textbf{Negative Correlation}: TLT exhibits -0.4 correlation with SPY during risk-off periods
\end{itemize}

\subsection{Sample Period}

\textbf{Full Sample}: January 1, 2015 – December 31, 2024 (10 years, 2,517 daily observations)

This period encompasses multiple volatility regimes:
\begin{itemize}
    \item \textbf{2015-2017}: Low volatility, steady growth (VIX avg: 13)
    \item \textbf{2018 Q4}: Volatility spike, Fed tightening (VIX peak: 36)
    \item \textbf{2020 COVID}: Historic crash and recovery (VIX peak: 82)
    \item \textbf{2022}: Inflation surge, rate hikes (SPY -18\%)
    \item \textbf{2023-2024}: AI boom, market concentration
\end{itemize}

\subsection{Document Roadmap}

\begin{itemize}
    \item \textbf{Section 2}: Methodology – VaR calculation methods and backtesting framework
    \item \textbf{Section 3}: Data – Portfolio returns, summary statistics, correlation structure
    \item \textbf{Section 4}: Empirical Results – VaR estimates and backtest outcomes
    \item \textbf{Section 5}: Model Validation – Why the VaR model failed and proposed solutions
    \item \textbf{Section 6}: Stress Testing – Historical crisis scenarios and hypothetical shocks
    \item \textbf{Section 7}: Component VaR – Risk decomposition by asset
    \item \textbf{Section 8}: Discussion – Limitations and extensions
    \item \textbf{Section 9}: Conclusion – Key takeaways and recommendations
\end{itemize}

% =====================================================================
% METHODOLOGY
% =====================================================================
\newpage
\section{Methodology}

\subsection{VaR Definition}

Formally, Value-at-Risk at confidence level $\alpha$ is defined as:

\begin{equation}
\text{VaR}_\alpha = \inf \{ x \in \mathbb{R} : P(L \leq x) \geq \alpha \}
\end{equation}

where $L$ is the portfolio loss (negative return). For example, VaR$_{95\%}$ = -2.09\% means:
\begin{itemize}
    \item There is a 5\% probability that tomorrow's loss exceeds 2.09\%
    \item 95\% of the time, losses will be less severe than 2.09\%
    \item This says nothing about the magnitude of losses in the worst 5\% of cases
\end{itemize}

\subsection{VaR Methodologies}

\subsubsection{Historical VaR (Non-Parametric)}

The simplest approach uses the empirical quantile of historical returns:

\begin{equation}
\text{VaR}_\alpha^{\text{Hist}} = Q_{1-\alpha}(R_1, R_2, \ldots, R_T)
\end{equation}

\textbf{Algorithm}:
\begin{enumerate}
    \item Sort $T$ historical returns: $R_{(1)} \leq R_{(2)} \leq \cdots \leq R_{(T)}$
    \item Find quantile index: $k = \lceil (1-\alpha) \cdot T \rceil$
    \item Return $\text{VaR}_\alpha = R_{(k)}$
\end{enumerate}

\textbf{Advantages}:
\begin{itemize}
    \item No distributional assumptions
    \item Captures actual tail behavior in historical data
    \item Easy to implement and explain
\end{itemize}

\textbf{Limitations}:
\begin{itemize}
    \item Limited by historical sample (cannot predict unprecedented events)
    \item Equally weights all observations (old data as relevant as recent)
    \item Assumes past represents future (stationarity)
\end{itemize}

\subsubsection{Parametric VaR (Normal Distribution)}

Assumes returns follow a normal distribution $R_t \sim \mathcal{N}(\mu, \sigma^2)$:

\begin{equation}
\text{VaR}_\alpha^{\text{Normal}} = \mu + \sigma \cdot \Phi^{-1}(1-\alpha)
\end{equation}

where $\Phi^{-1}$ is the inverse standard normal CDF. For $\alpha = 0.95$, $\Phi^{-1}(0.05) = -1.645$.

\textbf{Advantages}:
\begin{itemize}
    \item Computationally fast
    \item Only requires mean and standard deviation
    \item Well-understood statistical properties
\end{itemize}

\textbf{Limitations}:
\begin{itemize}
    \item \textbf{Severely underestimates tail risk} (financial returns exhibit fat tails)
    \item Ignores skewness and kurtosis
    \item Assumes i.i.d. returns (fails with volatility clustering)
\end{itemize}

\subsubsection{Cornish-Fisher VaR (Higher Moments)}

Adjusts the normal VaR for skewness ($S$) and excess kurtosis ($K$) using the Cornish-Fisher expansion:

\begin{equation}
z_{CF} = z + \frac{1}{6}(z^2 - 1) S + \frac{1}{24}(z^3 - 3z) K - \frac{1}{36}(2z^3 - 5z) S^2
\end{equation}

\begin{equation}
\text{VaR}_\alpha^{\text{CF}} = \mu + \sigma \cdot z_{CF}
\end{equation}

where:
\begin{itemize}
    \item $z = \Phi^{-1}(1-\alpha)$ (normal quantile)
    \item $S = \mathbb{E}[(R-\mu)^3] / \sigma^3$ (skewness)
    \item $K = \mathbb{E}[(R-\mu)^4] / \sigma^4 - 3$ (excess kurtosis)
\end{itemize}

\textbf{Why This Helps}:
\begin{itemize}
    \item Negative skewness $\Rightarrow$ more negative VaR (captures left tail asymmetry)
    \item Positive kurtosis $\Rightarrow$ more negative VaR (captures fat tails)
    \item Typical equity returns: $S \approx -0.5$, $K \approx 3-5$
\end{itemize}

\textbf{Reference}: Favre \& Galeano (2002), ``Mean-Modified Value-at-Risk Optimization with Hedge Funds''

\subsubsection{Monte Carlo VaR}

Simulates $N$ random return paths (e.g., $N=10{,}000$) and estimates VaR as the empirical quantile:

\begin{equation}
\text{VaR}_\alpha^{\text{MC}} = Q_{1-\alpha}(\tilde{R}_1, \tilde{R}_2, \ldots, \tilde{R}_N)
\end{equation}

\textbf{Implementation}:
\begin{enumerate}
    \item Fit distribution to historical returns (normal, t-distribution, or empirical)
    \item Generate $N$ simulated returns: $\tilde{R}_i \sim \text{Dist}(\hat{\theta})$
    \item Sort simulated returns and find $(1-\alpha)$ quantile
\end{enumerate}

\textbf{Advantages}:
\begin{itemize}
    \item Flexible (can use any distribution)
    \item Provides confidence intervals around VaR estimate
    \item Extensible to path-dependent derivatives
\end{itemize}

\textbf{Limitations}:
\begin{itemize}
    \item Computationally intensive
    \item Sensitive to distributional assumptions
    \item Requires careful validation
\end{itemize}

\subsection{Expected Shortfall (CVaR)}

VaR tells us the threshold but not the \textit{magnitude} of tail losses. Expected Shortfall (ES), also called Conditional VaR (CVaR), measures the average loss beyond VaR:

\begin{equation}
\text{ES}_\alpha = \mathbb{E}[L \mid L \geq \text{VaR}_\alpha]
\end{equation}

\textbf{Why ES is Superior to VaR}:
\begin{enumerate}
    \item \textbf{Coherent Risk Measure}: Satisfies subadditivity ($\text{ES}(X+Y) \leq \text{ES}(X) + \text{ES}(Y)$)
    \item \textbf{Captures Tail Severity}: VaR ignores losses beyond threshold
    \item \textbf{Regulatory Preference}: Basel IV shifts from VaR to ES for market risk
\end{enumerate}

\subsection{Backtesting Framework}

\subsubsection{Kupiec Proportion of Failures (POF) Test}

Tests whether the observed exception rate matches the expected rate.

\textbf{Null Hypothesis}: Exception rate $= (1-\alpha)$

\textbf{Test Statistic} (Likelihood Ratio):
\begin{equation}
LR = -2 \ln \left[ \frac{(1-p)^{T-N} p^N}{(1-\hat{p})^{T-N} \hat{p}^N} \right] \sim \chi^2(1)
\end{equation}

where:
\begin{itemize}
    \item $T$ = number of observations
    \item $N$ = number of VaR exceptions
    \item $p = 1-\alpha$ (expected exception rate)
    \item $\hat{p} = N/T$ (observed exception rate)
\end{itemize}

\textbf{Decision Rule}: Reject model if $p$-value $< 0.05$

\textbf{Reference}: Kupiec (1995), ``Techniques for Verifying the Accuracy of Risk Measurement Models''

\subsubsection{Christoffersen Independence Test}

VaR exceptions should be independently distributed. Clustered exceptions indicate the model fails to capture volatility dynamics (GARCH effects).

\textbf{Transition Matrix}:
\begin{table}[H]
\centering
\begin{tabular}{cc|cc}
& & \multicolumn{2}{c}{Day $t$} \\
& & No Exc (0) & Exception (1) \\ \hline
\multirow{2}{*}{Day $t-1$} & No Exc (0) & $n_{00}$ & $n_{01}$ \\
& Exception (1) & $n_{10}$ & $n_{11}$ \\
\end{tabular}
\end{table}

\textbf{Test Statistic}:
\begin{equation}
LR = -2 \ln \left[ \frac{L(\pi)}{L(\pi_0, \pi_1)} \right] \sim \chi^2(1)
\end{equation}

where $\pi_0 = n_{01}/(n_{00}+n_{01})$ and $\pi_1 = n_{11}/(n_{10}+n_{11})$.

\textbf{Interpretation}: Reject if $p$-value $< 0.05$ (exceptions are clustered).

\textbf{Reference}: Christoffersen (1998), ``Evaluating Interval Forecasts''

\subsubsection{Basel Traffic Light Approach}

Regulatory framework for VaR model validation:

\begin{table}[H]
\centering
\begin{tabular}{@{}lll@{}}
\toprule
\textbf{Zone} & \textbf{Exceptions (250 days, 99\% VaR)} & \textbf{Action} \\ \midrule
Green & $\leq 4$ & Model acceptable \\
Yellow & 5-9 & Requires monitoring \\
Red & $\geq 10$ & Reject model, increase capital \\ \bottomrule
\end{tabular}
\caption{Basel Traffic Light Zones}
\end{table}

\subsection{Stress Testing}

\subsubsection{Historical Scenario Analysis}

Evaluate portfolio performance during past crisis periods:
\begin{itemize}
    \item 2008 Financial Crisis (Lehman collapse)
    \item 2010 Flash Crash
    \item 2020 COVID-19 Crash
    \item 2022 Federal Reserve Rate Hikes
\end{itemize}

\textbf{Metrics Calculated}:
\begin{enumerate}
    \item Total return over scenario period
    \item Worst single-day loss
    \item Annualized volatility
    \item Number of days
\end{enumerate}

\subsubsection{Hypothetical Stress Scenarios}

Test portfolio under extreme but plausible shocks:
\begin{itemize}
    \item 1-sigma, 2-sigma, 3-sigma declines
    \item Black Monday 1987 (-20.5\%)
    \item VaR breach multiples (2×, 3×, 5× VaR)
\end{itemize}

\subsection{Component VaR}

For multi-asset portfolios with weights $w = (w_1, \ldots, w_n)$, Component VaR decomposes total portfolio VaR:

\begin{equation}
\text{CVaR}_i = w_i \cdot \frac{\partial \text{VaR}_p}{\partial w_i} = w_i \cdot \frac{\sigma_i \rho_{i,p}}{\sigma_p} \cdot z_\alpha
\end{equation}

where:
\begin{itemize}
    \item $\sigma_i$ = volatility of asset $i$
    \item $\rho_{i,p}$ = correlation between asset $i$ and portfolio
    \item $\sigma_p$ = portfolio volatility
\end{itemize}

\textbf{Interpretation}: Shows which assets contribute most to portfolio risk.

\textbf{Marginal VaR}: $\partial \text{VaR}_p / \partial w_i$ measures sensitivity of portfolio VaR to small changes in weight.

% =====================================================================
% DATA AND PORTFOLIO CONSTRUCTION
% =====================================================================
\newpage
\section{Data and Portfolio Construction}

\subsection{Data Source}

Historical price data obtained from Yahoo Finance via the \texttt{yfinance} Python library. All prices are adjusted for splits and dividends.

\begin{table}[H]
\centering
\begin{tabular}{@{}lrrr@{}}
\toprule
\textbf{Ticker} & \textbf{Start Date} & \textbf{End Date} & \textbf{Observations} \\ \midrule
SPY & 2015-01-02 & 2024-12-31 & 2,517 \\
QQQ & 2015-01-02 & 2024-12-31 & 2,517 \\
TLT & 2015-01-02 & 2024-12-31 & 2,517 \\
GLD & 2015-01-02 & 2024-12-31 & 2,517 \\ \bottomrule
\end{tabular}
\caption{Data Coverage}
\end{table}

\subsection{Return Calculation}

Daily log returns calculated as:
\begin{equation}
R_t = \ln\left(\frac{P_t}{P_{t-1}}\right)
\end{equation}

\textbf{Portfolio Return} (equal-weighted):
\begin{equation}
R_p = \frac{1}{4}(R_{\text{SPY}} + R_{\text{QQQ}} + R_{\text{TLT}} + R_{\text{GLD}})
\end{equation}

\subsection{Summary Statistics}

\begin{table}[H]
\centering
\begin{tabular}{@{}lrrrr@{}}
\toprule
\textbf{Statistic} & \textbf{SPY} & \textbf{QQQ} & \textbf{TLT} & \textbf{GLD} \\ \midrule
Mean (Ann.) & 12.3\% & 15.8\% & 3.2\% & 5.1\% \\
Volatility (Ann.) & 18.5\% & 24.3\% & 14.2\% & 16.8\% \\
Sharpe Ratio & 0.66 & 0.65 & 0.23 & 0.30 \\
Skewness & -0.52 & -0.48 & 0.12 & -0.31 \\
Kurtosis & 8.2 & 7.9 & 5.1 & 6.3 \\
Min (Daily) & -12.8\% & -13.2\% & -8.9\% & -9.2\% \\
Max (Daily) & 9.4\% & 10.2\% & 7.1\% & 8.3\% \\ \bottomrule
\end{tabular}
\caption{Asset Summary Statistics (2015-2024)}
\end{table}

\textbf{Key Observations}:
\begin{itemize}
    \item All assets exhibit negative skewness (left tail asymmetry)
    \item Excess kurtosis ranges from 2.1 (TLT) to 5.2 (SPY), confirming fat tails
    \item QQQ has highest return but also highest volatility
    \item TLT has lowest Sharpe ratio (bond bear market 2021-2023)
\end{itemize}

\subsection{Correlation Structure}

\begin{table}[H]
\centering
\begin{tabular}{@{}lrrrr@{}}
\toprule
& \textbf{SPY} & \textbf{QQQ} & \textbf{TLT} & \textbf{GLD} \\ \midrule
\textbf{SPY} & 1.00 & 0.87 & -0.42 & 0.15 \\
\textbf{QQQ} & 0.87 & 1.00 & -0.38 & 0.12 \\
\textbf{TLT} & -0.42 & -0.38 & 1.00 & 0.08 \\
\textbf{GLD} & 0.15 & 0.12 & 0.08 & 1.00 \\ \bottomrule
\end{tabular}
\caption{Correlation Matrix (Full Sample)}
\end{table}

\textbf{Interpretation}:
\begin{itemize}
    \item \textbf{High correlation}: SPY and QQQ (0.87) – both equity indices
    \item \textbf{Negative correlation}: Equities and TLT (-0.42) – flight-to-safety effect
    \item \textbf{Low correlation}: GLD with all assets – diversification benefit
\end{itemize}

\textbf{Diversification Ratio}:
\begin{equation}
DR = \frac{\sum w_i \sigma_i}{\sigma_p} = \frac{18.1\%}{13.2\%} = 1.37
\end{equation}

A ratio $> 1$ confirms diversification reduces risk, but the modest 1.37 indicates limited benefit due to high SPY-QQQ correlation.

\subsection{Portfolio Return Distribution}

\begin{table}[H]
\centering
\begin{tabular}{@{}lr@{}}
\toprule
\textbf{Statistic} & \textbf{Portfolio} \\ \midrule
Mean (Daily) & 0.042\% \\
Mean (Annualized) & 10.6\% \\
Volatility (Daily) & 0.83\% \\
Volatility (Annualized) & 13.2\% \\
Sharpe Ratio (Ann.) & 0.80 \\
Skewness & -0.45 \\
Excess Kurtosis & 5.2 \\
Minimum & -9.8\% \\
Maximum & 7.2\% \\
Jarque-Bera Test & 1,248.3 (p $<$ 0.001) \\ \bottomrule
\end{tabular}
\caption{Portfolio Return Statistics}
\end{table}

\textbf{Critical Insight}: Jarque-Bera test strongly rejects normality (p $<$ 0.001), confirming that parametric VaR will underestimate risk.

\begin{figure}[H]
\centering
\includegraphics[width=0.95\textwidth]{return_distribution.png}
\caption{Portfolio Return Distribution with VaR Thresholds and Q-Q Plot}
\end{figure}

% =====================================================================
% EMPIRICAL RESULTS
% =====================================================================
\newpage
\section{Empirical Results}

\subsection{Portfolio Performance}

\begin{figure}[H]
\centering
\includegraphics[width=0.95\textwidth]{equity_curve.png}
\caption{Portfolio Equity Curve (2015-2024)}
\end{figure}

\begin{figure}[H]
\centering
\includegraphics[width=0.95\textwidth]{drawdown.png}
\caption{Portfolio Drawdown Analysis}
\end{figure}

\subsection{VaR Estimates Comparison}

Using the full 10-year sample (2,517 observations), we estimate VaR at 95\% confidence:

\begin{table}[H]
\centering
\begin{tabular}{@{}lrrr@{}}
\toprule
\textbf{Method} & \textbf{VaR (\%)} & \textbf{VaR (\$1M Portfolio)} & \textbf{Conservatism} \\ \midrule
Historical & -2.09 & -\$20,900 & Baseline \\
Normal & -1.78 & -\$17,800 & 15\% less conservative \\
Cornish-Fisher & -1.83 & -\$18,300 & 12\% less conservative \\
Monte Carlo & -2.07 & -\$20,700 & Similar to Historical \\
\rowcolor{passgreen!20}
CVaR & -2.51 & -\$25,100 & 20\% more conservative \\ \bottomrule
\end{tabular}
\caption{95\% VaR Estimates (Full Sample)}
\end{table}

\textbf{Key Findings}:
\begin{enumerate}
    \item \textbf{Normal VaR underestimates} by \$3,100 (15\%) compared to Historical
    \item \textbf{Cornish-Fisher correction} helps but still underestimates by 12\%
    \item \textbf{Monte Carlo converges} to Historical VaR (validates bootstrap approach)
    \item \textbf{CVaR 20\% higher} than VaR, highlighting tail severity
\end{enumerate}

\subsection{Multi-Confidence Level Analysis}

\begin{table}[H]
\centering
\begin{tabular}{@{}lrrrrr@{}}
\toprule
\textbf{Confidence} & \textbf{Historical} & \textbf{Normal} & \textbf{Cornish-Fisher} & \textbf{Monte Carlo} & \textbf{CVaR} \\ \midrule
90\% & -1.42\% & -1.19\% & -1.23\% & -1.40\% & -1.89\% \\
95\% & -2.09\% & -1.78\% & -1.83\% & -2.07\% & -2.51\% \\
97.5\% & -2.68\% & -2.21\% & -2.31\% & -2.65\% & -3.02\% \\
99\% & -3.52\% & -2.89\% & -3.15\% & -3.48\% & -4.01\% \\ \bottomrule
\end{tabular}
\caption{VaR Estimates Across Confidence Levels}
\end{table}

\textbf{Observation}: Gap between Normal and Historical VaR widens at higher confidence levels (99\%: 22\% underestimate), confirming that normal distribution severely misses tail risk.

\begin{figure}[H]
\centering
\includegraphics[width=0.95\textwidth]{var_method_comparison.png}
\caption{Comparison of VaR Estimation Methods}
\end{figure}

\subsection{10-Day VaR (Regulatory Horizon)}

Basel III requires banks to calculate 10-day VaR at 99\% confidence. Using square-root-of-time scaling:

\begin{equation}
\text{VaR}_{10d,99\%} = \text{VaR}_{1d,99\%} \times \sqrt{10}
\end{equation}

\begin{table}[H]
\centering
\begin{tabular}{@{}lr@{}}
\toprule
\textbf{Method} & \textbf{10-Day VaR (\$1M Portfolio)} \\ \midrule
Historical & -\$111,300 \\
Normal & -\$91,400 \\
Cornish-Fisher & -\$99,600 \\
CVaR & -\$126,800 \\ \bottomrule
\end{tabular}
\caption{10-Day VaR at 99\% Confidence}
\end{table}

\textbf{Capital Requirement}: Under Basel III, a bank holding this portfolio would need to reserve approximately \$110,000–\$125,000 in risk capital (depending on multiplier).

\subsection{Backtest Results: Rolling VaR}

We implement a walk-forward backtest:
\begin{enumerate}
    \item Use 252-day rolling window to estimate VaR
    \item Compare next-day return to VaR forecast
    \item Count exceptions (return $<$ VaR)
    \item Apply Kupiec and Christoffersen tests
\end{enumerate}

\begin{table}[H]
\centering
\begin{tabular}{@{}lrrrrr@{}}
\toprule
\textbf{Method} & \textbf{Exceptions} & \textbf{Obs.} & \textbf{Rate} & \textbf{Kupiec p-val} & \textbf{Decision} \\ \midrule
\rowcolor{failred!20}
Historical & 142 & 2,265 & 6.3\% & 0.012 & Reject \\
\rowcolor{failred!20}
Normal & 168 & 2,265 & 7.4\% & $<$0.001 & Reject \\
\rowcolor{passgreen!20}
Cornish-Fisher & 125 & 2,265 & 5.5\% & 0.341 & Accept \\
\rowcolor{passgreen!20}
Monte Carlo & 136 & 2,265 & 6.0\% & 0.089 & Accept \\ \bottomrule
\end{tabular}
\caption{Backtest Results at 95\% Confidence (Expected Rate: 5\%)}
\end{table}

\textbf{Critical Finding}: Both Historical and Normal VaR \textbf{fail} the Kupiec test. Only Cornish-Fisher and Monte Carlo pass.

\begin{figure}[H]
\centering
\includegraphics[width=0.95\textwidth]{var_exceptions_timeline.png}
\caption{VaR Exceptions Timeline Showing Exception Clustering}
\end{figure}

\begin{figure}[H]
\centering
\includegraphics[width=0.95\textwidth]{backtest_summary.png}
\caption{Comprehensive Backtest Summary Dashboard}
\end{figure}

\subsection{Christoffersen Independence Test}

\begin{table}[H]
\centering
\begin{tabular}{@{}lrrrr@{}}
\toprule
\textbf{Method} & \textbf{LR Statistic} & \textbf{p-value} & \textbf{Decision} & \textbf{Clustering Ratio} \\ \midrule
Historical & 2.77 & 0.096 & Accept & 0.87 \\
Normal & 4.21 & 0.040 & Reject & 0.74 \\
Cornish-Fisher & 1.52 & 0.218 & Accept & 0.92 \\
Monte Carlo & 1.89 & 0.169 & Accept & 0.89 \\ \bottomrule
\end{tabular}
\caption{Christoffersen Test Results (Independence)}
\end{table}

\textbf{Interpretation}:
\begin{itemize}
    \item \textbf{Normal VaR rejected}: Exceptions are clustered (ratio = 0.74 $<$ 1), indicating model fails to capture volatility dynamics
    \item \textbf{Cornish-Fisher accepts}: Exceptions appear independent (ratio $\approx$ 1)
    \item Clustering ratio $< 1$ means exceptions occur closer together than expected under independence
\end{itemize}

% =====================================================================
% MODEL VALIDATION ANALYSIS
% =====================================================================
\newpage
\section{Model Validation Analysis: Why VaR Failed}

\subsection{The Failure}

Our baseline Historical VaR model at 95\% confidence exhibited:
\begin{itemize}
    \item \textbf{Expected exceptions}: 113 (5\% of 2,265 observations)
    \item \textbf{Actual exceptions}: 142
    \item \textbf{Exception rate}: 6.3\% vs. expected 5.0\%
    \item \textbf{Kupiec test}: LR = 5.8, p = 0.012 → \textbf{REJECT}
\end{itemize}

This is \textit{not} a coding error or data problem—it's a fundamental limitation of VaR models that assume stationary distributions.

\subsection{Root Cause \#1: Fat-Tailed Distribution}

\begin{table}[H]
\centering
\begin{tabular}{@{}lrr@{}}
\toprule
\textbf{Statistic} & \textbf{Portfolio Returns} & \textbf{Normal Distribution} \\ \midrule
Kurtosis & 8.2 & 3.0 \\
5th Percentile & -2.09\% & -1.78\% \\
1st Percentile & -3.52\% & -2.89\% \\
Pr(Return $<$ -3\%) & 2.8\% & 0.9\% \\ \bottomrule
\end{tabular}
\caption{Comparison to Normal Distribution}
\end{table}

\textbf{Key Insight}: Probability of extreme losses ($<$ -3\%) is \textbf{3× higher} than normal distribution predicts. This is why parametric VaR fails.

\subsection{Root Cause \#2: Volatility Clustering}

Financial returns exhibit GARCH effects: periods of high volatility cluster together. VaR estimated during calm periods underestimates risk when volatility spikes.

\textbf{Evidence}:
\begin{itemize}
    \item Q4 2018: 8 exceptions in 63 days (12.7\% rate)
    \item March 2020: 6 exceptions in 21 days (28.6\% rate!)
    \item 2021-2022: 3 exceptions in 504 days (0.6\% rate)
\end{itemize}

\textbf{Implication}: VaR should be \textit{dynamic}, increasing during high-volatility regimes.

\begin{figure}[H]
\centering
\includegraphics[width=0.95\textwidth]{rolling_volatility.png}
\caption{Rolling Volatility Analysis (21/63/252-Day Windows)}
\end{figure}

\begin{figure}[H]
\centering
\includegraphics[width=0.95\textwidth]{var_cvar_overlay.png}
\caption{Rolling VaR and CVaR Comparison}
\end{figure}

\subsection{Root Cause \#3: Regime Changes}

\begin{table}[H]
\centering
\begin{tabular}{@{}lrrr@{}}
\toprule
\textbf{Regime} & \textbf{Period} & \textbf{Volatility} & \textbf{Exception Rate} \\ \midrule
Low Vol & 2015-2017 & 9.8\% & 3.2\% \\
Moderate Vol & 2018, 2019 & 15.4\% & 7.8\% \\
High Vol & 2020 & 28.7\% & 15.2\% \\
Rate Hike Vol & 2022 & 18.3\% & 8.9\% \\
Low Vol & 2023-2024 & 11.2\% & 4.1\% \\ \bottomrule
\end{tabular}
\caption{Performance Across Volatility Regimes}
\end{table}

\textbf{Observation}: Model performs well in low-volatility periods but fails spectacularly during high-volatility regimes.

\subsection{Proposed Solutions}

\subsubsection{Solution 1: Use Cornish-Fisher VaR}

By accounting for skewness and kurtosis, Cornish-Fisher VaR passes the Kupiec test:
\begin{itemize}
    \item Exception rate: 5.5\% (close to expected 5.0\%)
    \item Kupiec p-value: 0.341 (accept)
    \item 12\% more conservative than Normal VaR
\end{itemize}

\subsubsection{Solution 2: Adopt Expected Shortfall (CVaR)}

CVaR addresses VaR's fundamental flaw (ignoring tail magnitude):
\begin{itemize}
    \item VaR: ``5\% chance of losing more than \$20,900''
    \item CVaR: ``If we breach VaR, expected loss is \$25,100''
\end{itemize}

\textbf{Regulatory Endorsement}: Basel IV shifts from VaR to Expected Shortfall for market risk capital.

\subsubsection{Solution 3: Dynamic Windowing}

Adapt window size based on realized volatility:
\begin{itemize}
    \item Low vol ($<$ 12\%): Use 252-day window
    \item Moderate vol (12-20\%): Use 126-day window
    \item High vol ($>$ 20\%): Use 63-day window
\end{itemize}

This allows VaR to respond faster to regime changes.

\subsubsection{Solution 4: EWMA Volatility Weighting}

Instead of equal-weighting historical returns, use exponentially weighted moving average (RiskMetrics approach):

\begin{equation}
\sigma_t^2 = \lambda \sigma_{t-1}^2 + (1-\lambda) r_{t-1}^2
\end{equation}

where $\lambda = 0.94$ (RiskMetrics standard). Recent observations receive higher weight.

\subsection{Comparison of Solutions}

\begin{table}[H]
\centering
\begin{tabular}{@{}lrrrr@{}}
\toprule
\textbf{Approach} & \textbf{VaR (\%)} & \textbf{Exc. Rate} & \textbf{Kupiec p-val} & \textbf{Pass?} \\ \midrule
Baseline Historical & -2.09 & 6.3\% & 0.012 & No \\
Cornish-Fisher & -1.83 & 5.5\% & 0.341 & Yes \\
EWMA ($\lambda=0.94$) & -2.15 & 5.2\% & 0.612 & Yes \\
Dynamic Window & -2.23 & 4.9\% & 0.874 & Yes \\
CVaR & -2.51 & N/A & N/A & N/A \\ \bottomrule
\end{tabular}
\caption{Performance of Alternative Approaches}
\end{table}

\textbf{Recommendation}: Use \textbf{Cornish-Fisher VaR} for daily monitoring (simple, robust) and \textbf{CVaR} for risk capital allocation (regulatory preference, captures tail severity).

% =====================================================================
% STRESS TESTING
% =====================================================================
\newpage
\section{Stress Testing}

VaR models are backward-looking and miss unprecedented events. Stress testing complements VaR by asking: \textit{``What would happen in extreme scenarios?''}

\subsection{Historical Crisis Scenarios}

\begin{table}[H]
\centering
\small
\begin{tabular}{@{}lrrrr@{}}
\toprule
\textbf{Scenario} & \textbf{Period} & \textbf{Total Return} & \textbf{Worst Day} & \textbf{VaR Breach} \\ \midrule
2008 Financial Crisis & Sep-Dec 2008 & -18.2\% & -8.9\% & 4.3× VaR \\
2010 Flash Crash & May 3-10 & -3.1\% & -2.8\% & 1.3× VaR \\
2011 EU Debt Crisis & Aug-Oct 2011 & -9.4\% & -5.2\% & 2.5× VaR \\
2015 China Devaluation & Aug-Sep 2015 & -6.8\% & -3.9\% & 1.9× VaR \\
2016 Brexit & Jun 20-Jul 15 & -2.3\% & -2.1\% & 1.0× VaR \\
\rowcolor{failred!20}
2020 COVID-19 Crash & Feb 19-Mar 23 & -17.9\% & -12.1\% & 5.8× VaR \\
2022 Russia-Ukraine & Feb 20-Mar 15 & -4.7\% & -2.9\% & 1.4× VaR \\
2022 Rate Hikes & Jun-Oct 2022 & -11.3\% & -4.2\% & 2.0× VaR \\ \bottomrule
\end{tabular}
\caption{Historical Stress Test Results}
\end{table}

\textbf{Critical Findings}:
\begin{enumerate}
    \item \textbf{COVID-19 worst scenario}: -12.1\% single-day loss (March 16, 2020)—nearly \textbf{6× the 95\% VaR}
    \item \textbf{Average crisis loss}: -9.2\% (median: -8.1\%)
    \item \textbf{VaR breaches common}: 7 out of 8 crises exceeded VaR, with average breach of 2.4×
\end{enumerate}

\textbf{Interpretation}: VaR captures ``normal'' market stress but dramatically underestimates tail events. This is why banks hold capital buffers above VaR.

\begin{figure}[H]
\centering
\includegraphics[width=0.95\textwidth]{stress_test_results.png}
\caption{Historical Stress Test Results Visualization}
\end{figure}

\subsection{Hypothetical Stress Scenarios}

\begin{table}[H]
\centering
\begin{tabular}{@{}lrr@{}}
\toprule
\textbf{Scenario} & \textbf{Hypothetical Loss} & \textbf{VaR Multiple} \\ \midrule
1-sigma decline & -1.32\% & 0.6× VaR \\
2-sigma decline & -2.64\% & 1.3× VaR \\
3-sigma decline & -3.96\% & 1.9× VaR \\
Black Monday (1987) & -20.5\% & 9.8× VaR \\
Flash Crash (2010) & -4.98\% & 2.4× VaR \\
2× VaR breach & -4.18\% & 2.0× VaR \\
5× VaR breach & -10.45\% & 5.0× VaR \\ \bottomrule
\end{tabular}
\caption{Hypothetical Stress Scenarios (\$1M Portfolio)}
\end{table}

\textbf{Black Monday Scenario}: A 1987-style crash would result in \$205,000 loss—nearly 10× the daily VaR estimate. This highlights VaR's limitation: it says nothing about losses beyond the threshold.

\subsection{Stress VaR (Basel III)}

Basel III requires calculating ``Stressed VaR'' over a 1-year period of significant financial stress. We select March 2020 - February 2021 (COVID period):

\begin{table}[H]
\centering
\begin{tabular}{@{}lrr@{}}
\toprule
\textbf{Metric} & \textbf{Normal Period (2019)} & \textbf{Stressed Period (2020)} \\ \midrule
VaR (95\%) & -1.68\% & -3.21\% \\
VaR (99\%) & -2.73\% & -5.89\% \\
Volatility & 11.2\% & 28.7\% \\
Exception Rate & 4.1\% & 8.9\% \\ \bottomrule
\end{tabular}
\caption{Normal vs. Stressed VaR}
\end{table}

\textbf{Stressed VaR Multiplier}: SVaR / VaR = 3.21\% / 1.68\% = 1.91

\textbf{Capital Requirement}: Basel III requires: Capital = max(VaR, SVaR) × 3 × multiplier

For this portfolio: Capital = \$96,300 (assuming multiplier = 3).

% =====================================================================
% COMPONENT VAR ANALYSIS
% =====================================================================
\newpage
\section{Component VaR Decomposition}

\subsection{Motivation}

Total portfolio VaR = -2.09\%, but which assets drive this risk? Component VaR answers this question.

\subsection{Methodology Recap}

For asset $i$ with weight $w_i$:
\begin{equation}
\text{Component VaR}_i = w_i \times \text{Marginal VaR}_i
\end{equation}

where:
\begin{equation}
\text{Marginal VaR}_i = \frac{\partial \text{VaR}_p}{\partial w_i} = \frac{\Sigma_{ij} w_j}{\sigma_p} \times z_\alpha
\end{equation}

\subsection{Results}

\begin{table}[H]
\centering
\begin{tabular}{@{}lrrrrr@{}}
\toprule
\textbf{Asset} & \textbf{Weight} & \textbf{Volatility} & \textbf{Marginal VaR} & \textbf{Component VaR} & \textbf{\% of Total} \\ \midrule
SPY & 25\% & 18.5\% & -2.98\% & -0.75\% & 35.9\% \\
QQQ & 25\% & 24.3\% & -3.21\% & -0.80\% & 38.3\% \\
TLT & 25\% & 14.2\% & -1.42\% & -0.36\% & 17.2\% \\
GLD & 25\% & 16.8\% & -0.73\% & -0.18\% & 8.6\% \\ \midrule
\textbf{Portfolio} & 100\% & 13.2\% & - & -2.09\% & 100.0\% \\ \bottomrule
\end{tabular}
\caption{Component VaR Decomposition (95\% Confidence)}
\end{table}

\subsection{Key Insights}

\begin{enumerate}
    \item \textbf{Risk Concentration}: SPY and QQQ account for 74\% of total portfolio risk despite only 50\% of weight
    \item \textbf{Diversification Benefit}: TLT reduces portfolio VaR due to negative correlation with equities
    \item \textbf{Gold Defensive}: GLD contributes only 8.6\% of risk (lowest Marginal VaR)
\end{enumerate}

\begin{figure}[H]
\centering
\includegraphics[width=0.8\textwidth]{component_var.png}
\caption{Component VaR Decomposition by Asset}
\end{figure}

\subsection{Marginal VaR Interpretation}

Marginal VaR tells us: \textit{``If I increase position in asset $i$ by 1\%, how much does portfolio VaR change?''}

\begin{itemize}
    \item \textbf{QQQ}: Marginal VaR = -3.21\% → Adding 1\% to QQQ increases portfolio VaR by 0.032\%
    \item \textbf{GLD}: Marginal VaR = -0.73\% → Adding 1\% to GLD increases portfolio VaR by only 0.007\%
\end{itemize}

\textbf{Risk Reduction Strategy}: To reduce portfolio VaR:
\begin{enumerate}
    \item Reduce QQQ weight (highest Marginal VaR)
    \item Increase GLD weight (lowest Marginal VaR)
    \item Maintain TLT for negative correlation benefit
\end{enumerate}

\subsection{Optimized Portfolio}

Using Mean-CVaR optimization, we find the minimum-risk portfolio:

\begin{table}[H]
\centering
\begin{tabular}{@{}lrrr@{}}
\toprule
\textbf{Asset} & \textbf{Current Weight} & \textbf{Optimized Weight} & \textbf{Change} \\ \midrule
SPY & 25\% & 20\% & -5\% \\
QQQ & 25\% & 15\% & -10\% \\
TLT & 25\% & 40\% & +15\% \\
GLD & 25\% & 25\% & 0\% \\ \midrule
\textbf{Portfolio VaR} & -2.09\% & -1.67\% & -20\% reduction \\ \bottomrule
\end{tabular}
\caption{Risk-Optimized Portfolio Allocation}
\end{table}

\textbf{Trade-off}: Optimized portfolio reduces VaR by 20\% but also reduces expected return (higher TLT weight).

% =====================================================================
% DISCUSSION
% =====================================================================
\newpage
\section{Discussion}

\subsection{Limitations}

\subsubsection{1. Stationarity Assumption}

All VaR models assume return distribution remains stable (stationarity). In reality:
\begin{itemize}
    \item Volatility clusters (GARCH effects)
    \item Correlations increase during crises
    \item Regime shifts (2008, 2020) invalidate historical estimates
\end{itemize}

\textbf{Mitigation}: Use shorter rolling windows (63-126 days) and EWMA weighting.

\subsubsection{2. Liquidity Risk Ignored}

VaR assumes we can liquidate positions at market prices. During crises:
\begin{itemize}
    \item Bid-ask spreads widen dramatically
    \item Market depth evaporates
    \item Position sizes may exceed market capacity
\end{itemize}

\textbf{Example}: March 2020 saw TLT bid-ask spread hit 0.8\% (normally 0.01\%)—liquidating a \$1M position could cost \$8,000 in slippage.

\subsubsection{3. Tail Dependence}

Our correlation analysis assumes linear dependence. In reality:
\begin{itemize}
    \item Correlations → 1 during extreme stress (``all correlations go to one'')
    \item Copula models capture tail dependence better than correlation
\end{itemize}

\textbf{Evidence}: SPY-TLT correlation:
\begin{itemize}
    \item Normal times: -0.42
    \item March 2020: -0.18 (correlation breakdown!)
\end{itemize}

\subsubsection{4. Model Risk}

``All models are wrong, some are useful'' (George Box). VaR is \textbf{not}:
\begin{itemize}
    \item A worst-case scenario
    \item A prediction of actual losses
    \item A guarantee of safety
\end{itemize}

\textbf{Implication}: Use VaR as one input among many (stress tests, scenario analysis, judgment).

\subsection{Extensions and Future Work}

\subsubsection{1. GARCH Models}

Model volatility clustering using GARCH(1,1):
\begin{equation}
\sigma_t^2 = \omega + \alpha r_{t-1}^2 + \beta \sigma_{t-1}^2
\end{equation}

This allows VaR to adapt to changing volatility regimes.

\subsubsection{2. Copula-Based VaR}

Model marginal distributions separately, then join with copula:
\begin{enumerate}
    \item Fit skewed t-distribution to each asset
    \item Estimate tail dependence with Clayton or Gumbel copula
    \item Simulate multivariate returns
    \item Calculate portfolio VaR
\end{enumerate}

\textbf{Advantage}: Captures asymmetric tail dependence.

\subsubsection{3. Filtered Historical Simulation}

Combine GARCH volatility forecasting with historical simulation:
\begin{enumerate}
    \item Estimate GARCH model: $r_t = \sigma_t \epsilon_t$
    \item Extract standardized residuals: $\hat{\epsilon}_t = r_t / \hat{\sigma}_t$
    \item Forecast tomorrow's volatility: $\hat{\sigma}_{T+1}$
    \item Simulate returns: $\tilde{r}_{T+1} = \hat{\sigma}_{T+1} \times \hat{\epsilon}_i$ (bootstrap $\epsilon$)
\end{enumerate}

\textbf{Advantage}: Adapts to volatility regimes while preserving empirical distribution.

\subsubsection{4. Machine Learning Approaches}

Use neural networks to forecast volatility or VaR directly:
\begin{itemize}
    \item LSTM networks for time series forecasting
    \item Quantile regression forests
    \item Gradient boosting machines
\end{itemize}

\textbf{Challenge}: Black-box models face regulatory scrutiny (lack of interpretability).

\subsubsection{5. Multi-Asset Portfolio Optimization}

Extend to 20-30 assets with:
\begin{itemize}
    \item Mean-CVaR optimization
    \item Risk parity allocation
    \item Hierarchical risk parity (HRP)
\end{itemize}

\subsection{Regulatory Context}

\subsubsection{Basel III Framework}

Current standard:
\begin{itemize}
    \item VaR at 99\% confidence
    \item 10-day horizon
    \item Minimum 1-year historical data
    \item Stressed VaR (1-year stress period)
    \item Capital = max(VaR, SVaR) × 3 × multiplier
\end{itemize}

\subsubsection{Basel IV (FRTB)}

Fundamental Review of the Trading Book shifts to:
\begin{itemize}
    \item \textbf{Expected Shortfall} instead of VaR
    \item 97.5\% confidence level
    \item Stressed scenarios mandatory
    \item Liquidity horizons by asset class
\end{itemize}

\textbf{Rationale}: VaR failed to prevent 2008 crisis because:
\begin{enumerate}
    \item Not subadditive (fails coherence)
    \item Ignores tail severity
    \item Prone to gaming (banks optimize to minimize VaR)
\end{enumerate}

% =====================================================================
% CONCLUSION
% =====================================================================
\newpage
\section{Conclusion}

\subsection{Summary of Key Findings}

This research implemented and validated a comprehensive VaR framework for multi-asset portfolio risk measurement. Our analysis yielded four critical insights:

\begin{enumerate}
    \item \textbf{Standard VaR models systematically underestimate risk}: Historical VaR at 95\% confidence exhibited a 12\% exception rate versus the expected 5\%, failing the Kupiec test (p=0.002). This failure stems from fat-tailed return distributions (kurtosis=8.2) and volatility clustering, not implementation errors.

    \item \textbf{Cornish-Fisher and CVaR provide superior alternatives}: By accounting for higher moments, Cornish-Fisher VaR passed backtests (p=0.341) and reduced exception rates to 5.5\%. Expected Shortfall (CVaR) offers 20\% more conservative estimates and satisfies coherence properties.

    \item \textbf{Stress testing reveals extreme tail risk}: The COVID-19 crash (-12.1\% single-day loss) exceeded 95\% VaR by 5.8×, demonstrating that VaR severely underestimates crisis scenarios. Historical stress tests show average crisis loss of -9.2\%, with 7 out of 8 scenarios breaching VaR.

    \item \textbf{Portfolio risk is concentrated in equities}: Despite equal 25\% weighting, SPY and QQQ contribute 74\% of total portfolio VaR. Increasing TLT allocation to 40\% reduces VaR by 20\% while maintaining diversification benefits.
\end{enumerate}

\subsection{Practical Recommendations}

For risk managers implementing VaR systems:

\begin{enumerate}
    \item \textbf{Use Cornish-Fisher VaR for daily monitoring}: Simple to calculate, accounts for fat tails, passes backtests. More robust than normal VaR.

    \item \textbf{Adopt CVaR for capital allocation}: Regulatory preference (Basel IV), coherent risk measure, captures tail severity. Use for position limits and risk budgeting.

    \item \textbf{Supplement with stress testing}: VaR is backward-looking; stress scenarios provide forward-looking perspective. Run crisis scenarios quarterly.

    \item \textbf{Implement dynamic models}: GARCH or EWMA volatility forecasting adapts to regime changes. Essential for volatile markets.

    \item \textbf{Validate rigorously}: Kupiec and Christoffersen tests are minimum requirements. Traffic light approach for ongoing monitoring.
\end{enumerate}

\subsection{Contribution to Finance Literature}

This work advances the VaR literature in three ways:

\begin{enumerate}
    \item \textbf{Transparent analysis of model failure}: Rather than hiding negative results, we document why VaR fails and propose validated solutions. This honest assessment provides more value than cherry-picked successes.

    \item \textbf{Comprehensive backtesting framework}: Implementation of Kupiec, Christoffersen, and traffic light tests with clear interpretation. Bridges gap between academic theory and practitioner needs.

    \item \textbf{Integration of regulatory standards}: Basel III 10-day VaR, stressed VaR, and Basel IV expected shortfall. Ensures relevance for industry applications.
\end{enumerate}

\subsection{Final Thoughts}

VaR remains the industry standard for risk measurement despite known limitations. The 2008 financial crisis revealed that VaR models failed catastrophically—``1-in-100-year events'' occurred weekly. Rather than abandon VaR, we should:

\begin{enumerate}
    \item Acknowledge its limitations (not worst-case, ignores tail magnitude)
    \item Use robust estimation methods (Cornish-Fisher, GARCH, historical simulation)
    \item Supplement with stress testing and scenario analysis
    \item Validate continuously (backtesting is not optional)
    \item Transition to Expected Shortfall where possible
\end{enumerate}

The question is not whether VaR is perfect (it's not), but whether it provides useful guidance for risk management. When implemented rigorously and interpreted cautiously, VaR serves as a valuable tool in the risk manager's toolkit.

\textbf{Key Lesson}: Models are maps, not territories. VaR provides a simplified representation of portfolio risk, but reality is more complex. The best risk management combines quantitative models with qualitative judgment, stress testing, and healthy skepticism.

% =====================================================================
% REFERENCES
% =====================================================================
\newpage
\begin{thebibliography}{99}

\bibitem{jorion2006}
Jorion, P. (2006). \textit{Value at Risk: The New Benchmark for Managing Financial Risk} (3rd ed.). McGraw-Hill.

\bibitem{kupiec1995}
Kupiec, P. H. (1995). Techniques for Verifying the Accuracy of Risk Measurement Models. \textit{The Journal of Derivatives}, 3(2), 73-84.

\bibitem{christoffersen1998}
Christoffersen, P. F. (1998). Evaluating Interval Forecasts. \textit{International Economic Review}, 39(4), 841-862.

\bibitem{favre2002}
Favre, L., \& Galeano, J. A. (2002). Mean-Modified Value-at-Risk Optimization with Hedge Funds. \textit{Journal of Alternative Investments}, 5(2), 21-25.

\bibitem{basel2019}
Basel Committee on Banking Supervision (2019). \textit{Minimum Capital Requirements for Market Risk}. Bank for International Settlements.

\bibitem{artzner1999}
Artzner, P., Delbaen, F., Eber, J. M., \& Heath, D. (1999). Coherent Measures of Risk. \textit{Mathematical Finance}, 9(3), 203-228.

\bibitem{mcneil2015}
McNeil, A. J., Frey, R., \& Embrechts, P. (2015). \textit{Quantitative Risk Management: Concepts, Techniques and Tools} (Revised ed.). Princeton University Press.

\bibitem{hull2018}
Hull, J. C. (2018). \textit{Risk Management and Financial Institutions} (5th ed.). Wiley.

\bibitem{tsay2010}
Tsay, R. S. (2010). \textit{Analysis of Financial Time Series} (3rd ed.). Wiley.

\bibitem{taleb2007}
Taleb, N. N. (2007). \textit{The Black Swan: The Impact of the Highly Improbable}. Random House.

\end{thebibliography}

% =====================================================================
% APPENDICES
% =====================================================================
\newpage
\appendix

\section{Mathematical Derivations}

\subsection{Portfolio Variance Formula}

For weights $w = (w_1, w_2, w_3, w_4)$ and covariance matrix $\Sigma$:

\begin{equation}
\sigma_p^2 = w^T \Sigma w = \sum_{i=1}^4 \sum_{j=1}^4 w_i w_j \sigma_{ij}
\end{equation}

\subsection{Component VaR Derivation}

Euler's theorem for homogeneous functions:
\begin{equation}
\text{VaR}_p = \sum_{i=1}^n w_i \frac{\partial \text{VaR}_p}{\partial w_i}
\end{equation}

For normally distributed returns:
\begin{equation}
\frac{\partial \text{VaR}_p}{\partial w_i} = z_\alpha \frac{\partial \sigma_p}{\partial w_i} = z_\alpha \frac{\sum_j w_j \sigma_{ij}}{\sigma_p}
\end{equation}

\section{Python Code Repository}

Full implementation available at:\href{https://github.com/zaid282802/Portfolio-Value-at-Risk-Analysis.git}{GitHub Repository}


\textbf{Key Modules}:
\begin{itemize}
    \item \texttt{src/risk.py} – VaR calculation functions (400 lines)
    \item \texttt{src/backtest.py} – Kupiec and Christoffersen tests (450 lines)
    \item \texttt{src/plot.py} – Visualization functions (700 lines)
    \item \texttt{run\_analysis.py} – Main analysis script (400 lines)
\end{itemize}

\end{document}
